\section{LIV-TS Methodology}
\label{sec:method}

LIV-TS has three components: (1) an \emph{operator pool} of LIV classes that form the search alphabet; (2) two \emph{structural templates} that specify how operator slots are arranged within a forecasting model; and (3) a \emph{search and training protocol} combining NSGA-II evolution with TiDAR-TS parallel training.
Figure~\ref{fig:overview} illustrates the overall pipeline.

\begin{figure*}[t]
\centering
\resizebox{\textwidth}{!}{%
\begin{tikzpicture}[
  font=\small,
  >=Stealth,
  stage/.style={draw=black!60, rounded corners=4pt, fill=#1,
                minimum width=3.4cm, minimum height=4.6cm},
  op/.style={draw, rounded corners=2pt, fill=#1,
             text width=2.6cm, minimum height=0.30cm,
             font=\tiny, align=center, inner sep=2pt},
  gc/.style={draw, fill=#1, minimum size=0.44cm,
             font=\tiny\bfseries, inner sep=0pt, align=center},
  arr/.style={->, very thick, draw=black!55},
  note/.style={font=\tiny, text=black!55, align=center, text width=2.6cm},
]

%% ---- (a) Operator Pool ----
\node[stage=blue!8]  at (0,0) {};
\node[font=\scriptsize\bfseries] at (0, 2.0) {(a) Operator Pool};
\node[op=blue!25]   at (0,  1.25) {\textsc{low\_rank}: SA-1\ldots4};
\node[op=green!25]  at (0,  0.75) {\textsc{semi\_sep}: CfC};
\node[op=orange!25] at (0,  0.25) {\textsc{toeplitz}: GConv-1,2};
\node[op=gray!25]   at (0, -0.25) {\textsc{diagonal}: GMemless};
\node[note] at (0, -1.1) {$|\mathcal{P}|{=}8$ base\\$+$ diff.\ variants\\genome alphabet $\mathcal{P}^N$};

%% Arrow a -> b
\draw[arr] (1.85,0) -- (2.65,0);

%% ---- (b) Structural Templates ----
\node[stage=yellow!12] at (4.5,0) {};
\node[font=\scriptsize\bfseries] at (4.5, 2.0) {(b) Structural Templates};

%% --- S-Mamba (left column, x=3.50) ---
\node[font=\tiny\bfseries, text=violet!80] at (3.50, 1.72) {S-Mamba};
%% Input
\node[draw=black!45, fill=gray!10, rounded corners=1pt,
      minimum width=1.18cm, minimum height=0.24cm, font=\tiny, inner sep=1pt,align=center]
  (smIn) at (3.50, 1.42) {In $(B,L,C)$};
%% Stage 1
\node[draw=violet!55, fill=violet!8, rounded corners=2pt,
      minimum width=1.18cm, minimum height=0.68cm] (smS1) at (3.50, 0.72) {};
\node[font=\tiny\bfseries, text=violet!80] at (3.50, 0.97) {Stage 1 (var.)};
\node[gc=green!40]  at (3.27, 0.68) {S};
\node[gc=blue!40]   at (3.73, 0.68) {L};
%% Stage 2
\node[draw=orange!55, fill=orange!8, rounded corners=2pt,
      minimum width=1.18cm, minimum height=0.68cm] (smS2) at (3.50, -0.06) {};
\node[font=\tiny\bfseries, text=orange!80!black] at (3.50, 0.19) {Stage 2 (tmp.)};
\node[gc=orange!40] at (3.27,-0.10) {T};
\node[gc=gray!40]   at (3.73,-0.10) {D};
%% Output
\node[draw=black!45, fill=gray!10, rounded corners=1pt,
      minimum width=1.18cm, minimum height=0.24cm, font=\tiny, inner sep=1pt, align=center]
  (smOut) at (3.50,-0.82) {Out $(B,H,C)$};
%% Flow arrows
\draw[->, thin, black!40] (smIn.south)  -- (smS1.north);
\draw[->, thin, black!40] (smS1.south)  -- (smS2.north);
\draw[->, thin, black!40] (smS2.south)  -- (smOut.north);

%% Divider between columns
\draw[black!18, dashed, very thin] (4.50,-1.72) -- (4.50, 1.88);

%% --- iTransformer (right column, x=5.50) ---
\node[font=\tiny\bfseries, text=teal!75] at (5.50, 1.72) {iTransformer};
%% Input
\node[draw=black!45, fill=gray!10, rounded corners=1pt,
      minimum width=1.18cm, minimum height=0.24cm, font=\tiny, inner sep=1pt, align=center]
  (itIn) at (5.50, 1.42) {In $(B,L,C)$};
%% Block 1
\node[draw=teal!55, fill=teal!8, rounded corners=2pt,
      minimum width=1.18cm, minimum height=0.68cm] (itB1) at (5.50, 0.72) {};
\node[font=\tiny\bfseries, text=teal!75] at (5.50, 0.97) {Block 1};
\node[gc=blue!40]  at (5.27, 0.74) {L};
\node[gc=gray!40]  at (5.73, 0.74) {D};
\node[font=\tiny, text=teal!60] at (5.27, 0.53) {cv};
\node[font=\tiny, text=teal!60] at (5.73, 0.53) {pv};
%% Block 2
\node[draw=teal!55, fill=teal!8, rounded corners=2pt,
      minimum width=1.18cm, minimum height=0.68cm] (itB2) at (5.50,-0.06) {};
\node[font=\tiny\bfseries, text=teal!75] at (5.50, 0.19) {Block 2};
\node[gc=orange!40] at (5.27,-0.04) {T};
\node[gc=gray!40]   at (5.73,-0.04) {D};
\node[font=\tiny, text=teal!60] at (5.27,-0.25) {cv};
\node[font=\tiny, text=teal!60] at (5.73,-0.25) {pv};
%% Output
\node[draw=black!45, fill=gray!10, rounded corners=1pt,
      minimum width=1.18cm, minimum height=0.24cm, font=\tiny, inner sep=1pt, align=center]
  (itOut) at (5.50,-0.82) {Out $(B,H,C)$};
%% Flow arrows
\draw[->, thin, black!40] (itIn.south)  -- (itB1.north);
\draw[->, thin, black!40] (itB1.south)  -- (itB2.north);
\draw[->, thin, black!40] (itB2.south)  -- (itOut.north);

%% Legend
\node[note] at (4.5,-1.40) {cv\,=\,cross-variate\\pv\,=\,per-variate\\S/L/T/D\,=\,SS/LR/TP/DG};

%% Arrow b -> c
\draw[arr] (6.35,0) -- (7.15,0);

%% ---- (c) NSGA-II ----
\node[stage=red!6] at (9.0,0) {};
\node[font=\scriptsize\bfseries] at (9.0, 2.0) {(c) NSGA-II Search};

%% Mini scatter: axes
\draw[->] (8.10,-1.55) -- (9.95,-1.55) node[right, font=\tiny] {Params $f_2$};
\draw[->] (8.10,-1.55) -- (8.10,-0.10) node[above, font=\tiny] {MSE $f_1$};
%% Dominated points (gray)
\foreach \px/\py in {1.0/0.7, 0.6/1.1, 1.3/0.9, 1.5/0.5, 0.8/0.4, 1.7/1.2} {
  \fill[black!25] ({8.10+\px*0.495},{-1.55+\py*0.605}) circle (1.8pt);
}
%% Pareto front (blue) with connecting line
\draw[blue!50, thin]
  ({8.10+0.20*0.495},{-1.55+1.40*0.605}) --
  ({8.10+0.45*0.495},{-1.55+1.05*0.605}) --
  ({8.10+0.75*0.495},{-1.55+0.72*0.605}) --
  ({8.10+1.10*0.495},{-1.55+0.45*0.605}) --
  ({8.10+1.50*0.495},{-1.55+0.25*0.605});
\foreach \px/\py in {0.20/1.40, 0.45/1.05, 0.75/0.72, 1.10/0.45, 1.50/0.25} {
  \fill[blue!80] ({8.10+\px*0.495},{-1.55+\py*0.605}) circle (2.5pt);
}
\node[font=\tiny, text=blue!70, align=left] at (9.65,-0.60) {Pareto\\front};

\node[font=\tiny, align=center] at (9.0, 0.90) {pop.\ 16,\ gen.\ 18};
\node[font=\tiny, align=center] at (9.0, 0.60) {500 steps / genome};
\node[font=\tiny, align=center] at (9.0, 0.30) {tournament select.};
\node[font=\tiny, align=center] at (9.0, 0.02) {crossover + mutation};
\node[font=\tiny, align=center] at (9.0,-0.26) {$p_{\mathrm{mut}}{=}0.10$};

%% Arrow c -> d
\draw[arr] (10.85,0) -- (11.65,0);

%% ---- (d) Retrain + Transfer ----
\node[stage=green!8] at (13.5,0) {};
\node[font=\scriptsize\bfseries] at (13.5, 2.0) {(d) Retrain \& Transfer};

\node[font=\tiny\bfseries]      at (13.5,  1.35) {Full retraining:};
\node[font=\tiny, align=center] at (13.5,  1.02) {top-$K$ Pareto genomes};
\node[font=\tiny, align=center] at (13.5,  0.72) {20{,}000 steps};

\draw[dashed, black!40] (12.75, 0.45) -- (14.25, 0.45);

\node[font=\tiny\bfseries]      at (13.5,  0.12) {Cross-dataset transfer:};
\node[font=\tiny] at (13.5, -0.18) {ETTh1 $\to$ ETTh2};
\node[font=\tiny] at (13.5, -0.44) {ETTh1 $\to$ Exchange};
\node[font=\tiny] at (13.5, -0.70) {no re-evolution};

\node[note] at (13.5, -1.25) {Pareto-optimal\\accuracy-compactness\\trade-off};

\end{tikzpicture}%
}
\caption{Overview of the LIV-TS pipeline.
  \textbf{(a)} Operator pool: 8 base LIV classes spanning four token-mixing families, plus differential variants.
  \textbf{(b)} Structural templates: S-Mamba splits the genome into a Stage~1 variate-mixing backbone and a Stage~2 temporal-refinement backbone; iTransformer interleaves cross-variate (cv) and per-variate (pv) sub-layers within each block. Colored cells are searchable LIV operator slots.
  \textbf{(c)} NSGA-II evolves the population over 18 generations; each candidate is evaluated for 500 steps and placed on the Pareto front of MSE ($f_1$) versus parameter count ($f_2$).
  \textbf{(d)} Top Pareto-front genomes are retrained for 20{,}000 steps and applied directly to transfer datasets without re-running evolution.}
\label{fig:overview}
\end{figure*}

\subsection{LIV-TS Operator Pool}
\label{sec:pool}

We select nine base operator classes spanning all four LIV token-mixing families; Table~\ref{tab:operator_pool} lists the full pool.
The attention variants (SA-1 through SA-4) use $T = CB^\top/\sqrt{d}$ with softmax and differ in key-value sharing: SA-1 is standard MHA, SA-2 adds a depthwise Conv1d pre-projection, SA-3 (MQA \cite{shazeer2019mqattn}) shares one KV head across all queries, and SA-4 (GQA \cite{ainslie2023gqa}) groups KV heads in pairs, interpolating between MHA and MQA.
The recurrence operators Rec-4 (CfC, $16{\times}$ expansion) and Rec-5 (Compact CfC, $2{\times}$ expansion) implement the complementary-gated SEMI\_SEPARABLE pattern from Section~\ref{sec:cfc_liv}.
GConv-1 uses a fixed 3-tap learned kernel; GConv-2 generates a 64-tap input-dependent kernel via a small MLP \cite{poli2023hyena}, providing long-range context at $\mathcal{O}(Ld)$ cost.
GMemless is a SwiGLU FFN \cite{shazeer2020glu} with $4{\times}$ internal expansion and no cross-token interaction.
Every base class has a \emph{differential} counterpart (IDs 10--18) that subtracts two independent instances of the same class, approximately doubling expressiveness at $2{\times}$ parameter cost.

\begin{table}[t]
\caption{LIV-TS operator pool. Expansion: ratio of internal to model dimension $d$. Token-mix abbreviations: LR = Low-Rank, SS = Semi-Sep., TP = Toeplitz, DG = Diagonal.}
\label{tab:operator_pool}
\centering
\footnotesize
\setlength{\tabcolsep}{4pt}
\begin{tabular}{clllc}
\toprule
ID & Name & Token Mix & Ch.\ Mix & Exp. \\
\midrule
1  & SA-1 (MHA)      & LR & Grouped  & $1\times$ \\
2  & SA-2 (Conv-MHA) & LR & Grouped  & $1\times$ \\
3  & SA-3 (MQA)      & LR & Grouped  & $1\times$ \\
4  & SA-4 (GQA)      & LR & Grouped  & $1\times$ \\
5  & Rec-4 (CfC)     & SS & Diagonal & $16\times$ \\
6  & GConv-1         & TP & Diagonal & $1\times$ \\
7  & GConv-2         & TP & Diagonal & $1\times$ \\
8  & GMemless (FFN)  & DG & Dense    & $4\times$ \\
9  & Rec-5 (Compact CfC) & SS & Diagonal & $2\times$ \\
\bottomrule
\multicolumn{5}{l}{Classes 10--18: differential variants (LIV$_1(x)-$LIV$_2(x)$) of classes 1--9.}
\end{tabular}
\end{table}

\subsection{Genome Encoding and Structural Templates}
\label{sec:genome}

A genome $g$ is a fixed-length list of integer operator class IDs, one per slot in the model.
The channel-mix type is derived deterministically from the token-mix family of the selected class (Table~\ref{tab:operator_pool}), reducing the search space to $|\mathcal{P}|^N$.

For \textbf{S-Mamba}, the genome splits evenly between a Stage~1 variate-mixing backbone (operators mix across $C$ variate tokens bidirectionally) and a Stage~2 temporal-refinement backbone (operators act within each $D$-dimensional variate embedding).
For \textbf{iTransformer}, each block interleaves a cross-variate (cv) and a per-variate (pv) sub-layer; genome entries $g_{2b}$ and $g_{2b+1}$ set the operators for block $b$.
Both structures share the $(B,L,C)\!\to\!(B,H,C)$ interface with RevIN \cite{kim2022revin} normalisation; Figure~\ref{fig:overview} panel~(b) illustrates the two template flows.

\subsection{TiDAR-TS Adaptation}
\label{sec:tidar_adapt}

We adapt TiDAR \cite{tidar2024} to continuous-valued regression with three changes.
\textbf{(1) MSE loss.} Both the AR and draft objectives use mean squared error on RevIN-restored outputs:
$\mathcal{L}_{\mathrm{AR}} = \mathrm{MSE}(\hat{x}_{2:L}, x_{2:L})$ and $\mathcal{L}_{\mathrm{Diff}} = \mathrm{MSE}(\hat{y}_{1:H}, y_{1:H})$.
\textbf{(2) Deterministic draft.} The draft head predicts the forecast mean directly; no stochastic denoising step is needed.
\textbf{(3) Fixed $\alpha$.} We set $\alpha = 1.0$ in Eq.~\eqref{eq:tidar_loss} for all experiments, providing moderate AR regularisation without dominating the parallel objective.
Figure~\ref{fig:tidar} illustrates the training setup and the three inference modes.

\begin{figure}[t]
\centering
\resizebox{0.97\columnwidth}{!}{%
\begin{tikzpicture}[
  font=\small, >=Stealth,
  lbl/.style={font=\tiny, text=black!60},
  arr/.style={->, draw=black!50, thin},
  tok/.style={draw=#1!65, fill=#1!18, rounded corners=1pt,
              minimum width=0.36cm, minimum height=0.30cm,
              font=\tiny\bfseries, inner sep=0pt, align=center},
  blk/.style={draw=#1!65, fill=#1!12, rounded corners=3pt,
              minimum width=3.0cm, minimum height=0.44cm,
              font=\scriptsize, align=center},
]

%% ====================================================
%% (a) Training: tokens -> backbone -> dual heads
%% ====================================================
\node[font=\small\bfseries] at (1.55, 4.55) {(a) Training};

%% Token row
\foreach \i in {0,...,3} {
  \node[tok=blue] at ({0.05+\i*0.38}, 4.10) {};
}
\node[lbl,font=\tiny\bfseries,text=blue!75] at (0.62, 3.88) {lookback ($L$)};
\draw[black!25, thin, dashed] (1.64, 3.76) -- (1.64, 4.28);
\foreach \i in {0,...,2} {
  \node[tok=orange] at ({1.80+\i*0.38}, 4.10) {};
}
\node[lbl,font=\tiny\bfseries,text=orange!80!black] at (2.15, 3.88) {horizon ($H$)};

%% Backbone
\node[blk=black, minimum width=3.1cm] (bb) at (1.55, 3.28) {Backbone (hybrid mask)};
\foreach \xv in {0.05,0.43,0.81,1.19,1.80,2.18,2.56} {
  \draw[arr] (\xv, 3.95) -- (\xv, 3.51);
}

%% Two heads
\node[blk=blue, minimum width=1.4cm] (arh) at (0.72, 2.55) {AR Head};
\node[blk=orange, minimum width=1.4cm] (dfh) at (2.45, 2.55) {Draft Head};
\draw[arr] (bb.south -| arh) -- (arh.north);
\draw[arr] (bb.south -| dfh) -- (dfh.north);

%% Loss labels
\node[lbl,text=blue!70,align=center] at (0.72, 2.08) {$\mathcal{L}_{\mathrm{AR}}$\\{\tiny next-step on $x_{2..L}$}};
\node[lbl,text=orange!75,align=center] at (2.45, 2.08) {$\mathcal{L}_{\mathrm{Diff}}$\\{\tiny parallel on $y_{1..H}$}};
\draw[arr] (arh.south) -- (0.72, 1.80);
\draw[arr] (dfh.south) -- (2.45, 1.80);

%% Combined loss
\draw[arr] (0.72, 1.64) |- (1.55, 1.46);
\draw[arr] (2.45, 1.64) |- (1.55, 1.46);
\node[blk=gray, minimum width=3.3cm] at (1.55, 1.24) {$\mathcal{L}=\dfrac{\alpha\,\mathcal{L}_{\mathrm{AR}}+\mathcal{L}_{\mathrm{Diff}}}{1+\alpha}$};

%% Divider
\draw[black!18, dashed, thin] (3.48, 0.72) -- (3.48, 4.65);

%% ====================================================
%% (b) Inference modes
%% ====================================================
\node[font=\small\bfseries] at (6.10, 4.55) {(b) Inference Modes};

%% ---- Row 1: Pure draft ----
\node[lbl,font=\scriptsize\bfseries,text=green!60!black,anchor=east] at (4.22,4.10) {Draft};
\node[lbl,anchor=east] at (4.22,3.90) {$\mathcal{O}(1)$};
\foreach \i in {0,...,3} {
  \node[tok=blue] at ({4.38+\i*0.36},4.10) {};
}
\draw[arr] (5.80,4.10) -- (6.05,4.10);
\node[lbl,font=\tiny,align=center] at (6.35,4.10) {1 fwd};
\draw[arr] (6.65,4.10) -- (6.88,4.10);
\foreach \i in {0,...,3} {
  \node[tok=green] at ({6.99+\i*0.36},4.10) {};
}
\node[lbl,font=\tiny,text=green!60!black,anchor=west] at (8.45,4.10) {$\times 1$};

%% ---- Row 2: Partial-AR ----
\node[lbl,font=\scriptsize\bfseries,text=orange!70!black,anchor=east] at (4.22,3.20) {Partial};
\node[lbl,anchor=east] at (4.22,3.00) {$\mathcal{O}(k)$};
\foreach \i in {0,...,3} {
  \node[tok=blue] at ({4.38+\i*0.36},3.20) {};
}
\draw[arr] (5.80,3.20) -- (6.05,3.20);
\node[lbl,font=\tiny,align=center] at (6.35,3.20) {$k{+}1$ fwd};
\draw[arr] (6.65,3.20) -- (6.88,3.20);
\foreach \i in {0,1} {
  \node[tok=orange] at ({6.99+\i*0.36},3.20) {};
}
\foreach \i in {2,3} {
  \node[tok=green] at ({6.99+\i*0.36},3.20) {};
}
\node[lbl,font=\tiny,text=orange!70,anchor=west] at (6.96,2.98) {$\underbrace{\rule{0.72cm}{0pt}}_{k}$};
\node[lbl,font=\tiny,text=green!60!black,anchor=west] at (7.72,2.98) {$\underbrace{\rule{0.72cm}{0pt}}_{H-k}$};

%% ---- Row 3: Full-AR ----
\node[lbl,font=\scriptsize\bfseries,text=orange!80!black,anchor=east] at (4.22,2.30) {Full-AR};
\node[lbl,anchor=east] at (4.22,2.10) {$\mathcal{O}(H)$};
\foreach \i in {0,...,3} {
  \node[tok=blue] at ({4.38+\i*0.36},2.30) {};
}
\draw[arr] (5.80,2.30) -- (6.05,2.30);
\node[lbl,font=\tiny,align=center] at (6.35,2.30) {$H$ fwd};
\draw[arr] (6.65,2.30) -- (6.88,2.30);
\foreach \i in {0,...,3} {
  \node[tok=orange] at ({6.99+\i*0.36},2.30) {};
}
\node[lbl,font=\tiny,text=orange!80,anchor=west] at (8.45,2.30) {$\times H$};

%% Legend
\node[tok=blue]   at (4.38,1.38) {}; \node[lbl,anchor=west] at (4.60,1.38) {lookback};
\node[tok=green]  at (5.68,1.38) {}; \node[lbl,anchor=west] at (5.90,1.38) {draft (parallel)};
\node[tok=orange] at (7.26,1.38) {}; \node[lbl,anchor=west] at (7.48,1.38) {AR (sequential)};

\end{tikzpicture}%
}
\caption{TiDAR-TS overview.
  \textbf{(a) Training}: the backbone processes $L{+}H$ tokens under a hybrid mask; an AR head predicts $\mathcal{L}_{\mathrm{AR}}$ on the lookback and a Draft head predicts $\mathcal{L}_{\mathrm{Diff}}$ on all $H$ horizon steps in parallel.
  \textbf{(b) Inference modes}: \emph{Draft} ($\mathcal{O}(1)$) predicts all $H$ steps in one pass; \emph{Partial-AR} ($\mathcal{O}(k)$) refines the first $k$ steps; \emph{Full-AR} ($\mathcal{O}(H)$) applies $H$ sequential passes.}
\label{fig:tidar}
\end{figure}

\subsection{NSGA-II Search Protocol}
\label{sec:search}

We minimise two objectives jointly: validation MSE $f_1(g) = \mathcal{L}_{\mathrm{val}}(g)$ after $N_{\mathrm{eval}}$ training steps, and total parameter count $f_2(g) = \Phi(g)$.
No weighting is applied; NSGA-II tracks the Pareto front of accuracy versus compactness automatically.
We run 18 generations with population size 16; each candidate genome is evaluated for 500 Adam steps ($\mathrm{lr}=10^{-3}$) and its objectives recorded.
Offspring are produced by multi-point crossover and per-gene mutation ($p_{\mathrm{mut}}=0.10$); 2 elite individuals carry over unchanged each generation.
After search, the top-$K$ Pareto-front genomes (ranked by validation MSE) are retrained from scratch for 20{,}000 steps to obtain final test performance.
